%-*-latex-*-
%-*-latex-*-
\newcommand\COURSE{ciss245}
\newcommand\ASSESSMENT{q6401}
\newcommand\ASSESSMENTTYPE{Quiz}
\newcommand\POINTS{\textwhite{xxx/xxx}}

\input{myquizpreamble}
\input{yliow}
\input{\COURSE}
%-*-latex-*-

%-*-latex-*-

%-*-latex-*-
\renewcommand\TITLE{\ASSESSMENTTYPE \ \ASSESSMENT}

\renewcommand\EMAIL{}
\renewcommand\AUTHOR{eagraham1@cougars.ccis.edu}

\textwidth=6in
\begin{document}
\topmatterthree

%------------------------------------------------------------------------------
\nextq
In the Wumpus world, suppose you have several Wumpuses.
When a hunter kills a Wumpus, say \verb!joe_wumpus!, the hunter earns
\verb!joe_wumpus.points_!.
You think it's reasonable for
\verb!joe_wumpus.points_! to be 1000.
When the hunter kills any Wumpus, the number of points
awarded is the same.
Of course each Wumpus has a row and a column for its position in the
Wumpus World.
Each Wumpus also has a direction (N, S, E, W) which is the direction
that it is facing, i.e., if the Wumpus is facing N, then when it moves
it will move in the north direction if it does not bump against a wall.

Correct the following header file to optimize on memory usage
when you have a huge number of Wumpuses.
\\
\ANSWER
\begin{answercode}
// file: Wumpus.h
class Wumpus
{
public:
    Wumpus()
    : is_alive_(true), row_(0), col_(0), direction_('N')
    {}
    
    int get_points() const;
    // other methods not shown
private:
    bool is_alive_;
    int row_;
    int col_;
    char direction_;
    static const int points_ = 1000;
};
\end{answercode}

%------------------------------------------------------------------------------
\nextq
Refer to the question above.
Complete the following by implementing
\verb!Wumpus::Wumpus()!,
\verb!Wumpus::get_points()!, and whatever is necessary.
You need not implement methods not shown in the previous question.
(There is something else besides the above two functions.)
\\
\ANSWER
\begin{answercode}
// file: Wumpus.cpp
#include "Wumpus.h"
int get_points() const
{
  return points;
}

\end{answercode}

%------------------------------------------------------------------------------
\nextq
The \verb!DollarAmount! class behaves like this:
\begin{Verbatim}[frame=single,fontsize=\footnotesize]
DollarAmount amt0(1);
DollarAmount amt1(1234);
std::cout << amt0 << '\n'; // prints "$0.01" (without double-quotes)
std::cout << amt1 << '\n'; // prints "$12.34" (without double-quotes)
DollarAmount::set_format(1);
std::cout << amt0 << '\n'; // prints "1c" (without double-quotes)
std::cout << amt1 << '\n'; // prints "1234c" (without double-quotes)
DollarAmount::set_format(0);
std::cout << amt0 << '\n'; // prints "$0.01" (without double-quotes)
std::cout << amt1 << '\n'; // prints "$12.34" (without double-quotes)
\end{Verbatim}
Complete and correct \verb!DollarAmount.h! so that the
above code fragment works.
Do not define the bodies of the methods (do that in the next question).
\\
\ANSWER
\begin{answercode}
// file: DollarAmount.h

#ifndef DOLLARAMOUNT_H
#define DOLLARAMOUNT_H

#include <iostream>

class DollarAmount
{
public:
DollarAmount(int cents);
static void set_format(int format);
private:
    int cents_;
    int format_;
};

std::ostream& operator<<(std::ostream& cout, const DollarAmount& amt);
\end{answercode}

%------------------------------------------------------------------------------
\nextq
Refer to the previous question.
Complete \verb!DollarAmount.cpp! so that the code fragment in the previous
question works.
\\
\ANSWER
\begin{answercode}
// file: DollarAmount.cpp

#include <iostream>
#include "DollarAmount.h"
DollarAmount::DollarAmount(int cents)
: cents_(cents)
{}
void DollarAmount::set_format(int format)
{
  format_ = f;
}

std::ostream& operator<<(std::ostream& cout, const DollarAmount)
{
  if(DollarAmount::format__ == 0)
  {
    int dollars = amt.cents_ / 100;
    int cents = amt.cent_ % 100;
    cout << '$' << dollars << std::setw(2) << std::setfill('0') << cents;
  }
  else
  {
    cout << amt.cents_ << 'c';
  }
  
  return cout;
}
\end{answercode}

%------------------------------------------------------------------------------
%-*-latex-*-
\newpage
\input{instructions}
\end{document}

